\vspace{-5pt}
\section{Discussion}
\label{sec:discussion}

\parabf{Scope of the resource-profile model.}
\sysname models each function with a stable CPU-to-memory demand ratio and a
variable demand magnitude. This model is broadly applicable to fine-grained
serverless functions because many of them are single-purpose and exhibit stable
CPU-to-peak-memory profiles, as shown in
Figure~\ref{fig:ratio_stability_scatter}.

Single-resource-dominant streaming functions, such as video-processing and
uploader in Figure~\ref{fig:ratio_stability_scatter}, can still be represented
within the current model by assigning a conservative fixed resource ratio. This
may over-allocate the less variable resource but preserves normal system operation.
Providing specialized support for streaming functions would require extending
both the demand model and the SwapCapacity routine to allow a function's demand
ratio to vary with request size.

\looseness=-1
The current model does not directly support complex functions whose resource
consumption pattern changes qualitatively across requests. For such cases, a
natural extension is to classify requests into several virtual functions,
each with a single resource profile, and then apply the existing routing model to
those virtual functions. A more fundamental approach is to refactor such
functions by splitting branching or multi-stage processing logic into more
independent tasks, making their resource behavior easier for \sysname or other resource schedulers
to model stably. We leave these extensions to future work.

\parabf{Vertical scaling versus size-specific pools.}
The central challenge in predictive horizontal scaling with size-specific
prewarmed pools is determining function-specific boundaries that partition
requests into resource-size classes. Coarse-grained classes retain substantial
within-class demand variation and therefore require over-provisioning, whereas
fine-grained classes fragment warm capacity across pools and reduce instance-reuse
opportunities. Monotonic vertical scaling avoids this granularity trade-off by
allowing a warm instance to serve requests with different resource demands
without predefined size boundaries. We acknowledge that size-specific pools can
be effective when a function serves only a few stable request types, or when its
invocation volume is sufficiently high to maintain high reuse within every
fine-grained pool. Outside these regimes, vertical scaling provides a more
flexible way to accommodate heterogeneous request demands.
